% 原書 pp.347--352（添付PDF pp.369--374、文献一覧直前まで）
\sectionlink*{参考文献}
\label{sec:references-guide}

この2巻で扱った主題に関する文献は、優に5000編を超え、おそらくその倍近くに達する。このように多くの分野で研究したすべての人々の貢献を正確に伝えることは、本書の範囲を超える仕事である。そこで、個々の主題をさらに学び、文献を読み始める際に役立つ程度に短い文献一覧をまとめようとした。その結果、多くの総説に加え、本書の扱いを補うために必要だと感じた原論文を挙げることになった。

「e-print: hep-th/yymmnnn」と記した論文は、Los Alamosの物理学電子プレプリント・アーカイブから電子的に入手できる。
\begin{center}
\url{http://xxx.lanl.gov/abs/hep-th/yymmnnn}
\end{center}
訂正一覧およびその他の有用な情報については、次を参照されたい。
\begin{center}
\url{http://www.itp.ucsb.edu/~joep/bigbook.html}
\end{center}

\subsection*{一般的な参考文献}
読者には、Green, Schwarz, \& Witten (1987)\refcite{46}の2巻本（以下GSW）が有用な補助文献になるだろう。この本で用いられている共変的手法と光円錐の手法の組合せ、および重点を置いている発展的な主題は、本書の進め方をある程度補完している。Polyakov (1987)\refcite{86}の本は、共形場理論と弦の経路積分による扱いを重視しているが、その焦点は統一弦理論よりも、強い相互作用の弦理論とランダムな曲面の統計力学に置かれている。弦理論についてのほかの講義として、Peskin (1987)\refcite{79}、Lüst \& Theisen (1989)\refcite{63}、Alvarez-Gaumé \& Vazquez-Mozo (1992)\refcite{4}、D'Hoker (1993)\refcite{22}、Ooguri \& Yin (1997)\refcite{78}、Kiritsis (1997)\refcite{56}がある。

\subsection*{略史}
弦理論の歴史には、研究が非常に活発になった時期が3度あり、おおよそ1968--73年、1984--89年、そして1995年から現在までに当たる。最初の時期は、強い相互作用の$S$行列に対する現象論的な提案としてのVeneziano振幅から始まった。これが双対共鳴模型につながり、そこではVenezianoの公式が一般化され、体系化された。続いて、双対模型と弦との関係、臨界次元（26）、超共形不変性、そして臨界次元10の弦の存在が発見された。この最初の活発な時期は、強い相互作用の正しい理論としてQCDが発見されたことで終わった。しかし、弦理論に取り組んだ人々には、自分たちは豊かな物理理論を発見したのであり、自然界の何かに適用されるはずだという信念が残された。

その後の「静かな」時期には、少数の理論家が重力とゲージ相互作用の統一理論として弦理論を発展させた。この時期に時空の超対称性を持つ弦理論が発見され、超弦が量子重力の有限な理論であるとの議論がなされた。1984年の3つの発見により、弦理論が標準模型の統一にとっても有力な候補であることが確立した。これらの発見は、10次元におけるアノマリーのないカイラルな理論の存在、ヘテロ弦、そしてCalabi--Yauコンパクト化であった。その後の活発な時期には、主としてこれらの帰結を明らかにすることに力が注がれた。当時残され、現在も残っている最大の難点は、可能な真空の数が膨大であることと、その中から真空を選ぶ力学についての理解が不十分であることである。

最新の活発な時期は、新しい力学的な道具によって特徴づけられる。やはり、多くの重要な結果はその前の静かな時期にまでさかのぼるが、その重要性と相互の関係が明らかになるには時間を要した。鍵となった段階は、まずゲージ理論で、次に弦理論で、超対称性の帰結を体系的に解析したことであった。これによって多くの新しい発見が生まれた。弦理論における弱・強結合双対性、すべての弦理論が単一の理論の極限として統一されること、新しい第11の次元が現れる極限、Dブレーンの意義、そしてブラックホールの量子力学についての新しい理解である。

1984年以降の発展は、以下の文献と総説および第II巻で扱われているが、それ以前の時期については扱いがやや手薄になっている。1985年までの主要な論文の多くはSchwarz (1985)\refcite{94}に収録されており、そこには歴史の概説も含まれている。GSWの参考文献の部の第1章と第2章には、豊富な参考文献とともに、双対模型の時期についてさらに詳しい歴史が記されている。Scherk (1975)\refcite{90}の総説もこの時期を扱っている。

\subsection*{第1章}
現在では標準模型についての解説は数多くある。そのうち2つだけ挙げると、Quigg (1983)\refcite{87}は概観を、Cheng \& Li (1991)\refcite{19}はより詳細な解説を与えている。Weinberg (1979)\refcite{109}は、場の理論における量子重力の紫外問題を論じている。さまざまな統一の考え方についての参考文献は、それぞれに対応する章で挙げる。

弦の古典運動と光円錐量子化は、Goddard, Goldstone, Rebbi, \& Thorn (1973)\refcite{43}、およびRebbi (1974)\refcite{88}とScherk (1975)\refcite{90}による総説で論じられている。最初の2つにはLorentz不変性の丁寧な扱いが含まれる。本書で与えた発見的な議論はBrink \& Nielsen (1973)\refcite{11}に基づく。Weinberg (1964)\refcite{108}は、質量のないベクトル粒子とテンソル粒子の結合を論じている。

\subsection*{第2章}
一般的な参考文献に加えて、Friedan, Martinec, \& Shenker (1986)\refcite{37}は、本章の主題の多くを扱っている。Belavin, Polyakov, \& Zamolodchikov (1984)\refcite{9}はCFTの古典的な文献だが、第15章で扱う発展的な主題に向けられている。Ginsparg (1990)\refcite{40}はCFTの総説である。Gatto (1973)\refcite{38}とFradkin \& Palchik (1978)\refcite{32}は、4次元の量子場理論における共形不変性を扱っている。これは2次元の対称性よりはるかに小さいものの、多くの同じ考え方が適用できる。

\subsection*{第3章}
Polyakov経路積分についての追加の参考文献は、Polyakov (1981)\refcite{85}、Friedan (1982)\refcite{35}、Alvarez (1983)\refcite{1}である。Faddeev--Popovの手続きは、現代的な場の理論の教科書すべてで扱われている。基本的な手順は原論文Faddeev \& Popov (1967)\refcite{31}に従う。Birrell \& Davies (1982)\refcite{10}の本は、さまざまな次元でのWeylアノマリーを論じている（ただし、代わりに一般的な用語である共形アノマリーを用いている）。重力アノマリーはAlvarez-Gaumé \& Witten (1983)\refcite{5}で論じられている。

光円錐ゲージにおける頂点演算子のLSZに似た導出はMandelstam (1973)\refcite{64}にある。GSW (1987)\refcite{46}の第11章も参照されたい。一般の非線形シグマ模型のベータ関数の計算については、Alvarez-Gaumé, Freedman, \& Mukhi (1981)\refcite{2}とFriedan (1985)\refcite{36}を参照されたい。背景場中の弦についての一般的な参考文献はCallan, Friedan, Martinec, \& Perry (1985)\refcite{15}であり、総説にはCallan \& Thorlacius (1989)\refcite{17}とTseytlin (1989)\refcite{105}がある。線形ディラトン理論は、Chodos \& Thorn (1974)\refcite{20}ではCFTとして、Myers (1987)\refcite{74}では弦の背景として論じられている。

\subsection*{第4章}
OCQのノーゴースト定理の古典的な証明は、Brower (1972)\refcite{12}とGoddard \& Thorn (1972)\refcite{44}にある。拘束のある系の量子化に関する基本文献はDirac (1950)\refcite{27}であり、総説にはHenneaux (1985)\refcite{52}がある。ゲージ理論のBRST対称性は現代的な場の理論の教科書で論じられている。そのコホモロジーの構造はKugo \& Ojima (1979)\refcite{58}で論じられている。場の理論のBRST演算子はKato \& Ogawa (1983)\refcite{53}で導出されており、本書のノーゴースト定理の証明も、その論文の証明を大まかな手本としている。BRSTの枠組みでのほかの証明は、Freeman \& Olive (1986)\refcite{33}、Frenkel, Garland, \& Zuckerman (1986)\refcite{34}、Thorn (1987)\refcite{103}にある。

\subsection*{第5章}
D'Hoker \& Phong (1988)\refcite{25}は、この章と後続の章の主題についての広範な総説である。モジュライ空間上の測度はMoore \& Nelson (1986)\refcite{72}とD'Hoker \& Phong (1986)\refcite{24}で求められている。Alvarez (1983)\refcite{1}も参照されたい。遷移関数を用いた測度の表現は、Martinec (1987)\refcite{68}とAlvarez-Gaumé, Gomez, Moore, \& Vafa (1988)\refcite{3}で論じられている。振幅の一般的な性質は、Friedan, Martinec, \& Shenker (1986)\refcite{37}とMansfield (1987)\refcite{66}でも論じられている。

\subsection*{第6章}
樹木レベルの振幅の物理の一部はVeneziano (1974)\refcite{107}で論じられている。長距離極限でゲージ相互作用と重力相互作用が現れることは、Neveu \& Scherk (1972)\refcite{77}、Yoneya (1974)\refcite{112}、Scherk \& Schwarz (1974)\refcite{91}で展開されている。開弦のゲージ群に対する制約は、Marcus \& Sagnotti (1982)\refcite{67}とSchwarz (1982)\refcite{93}で論じられている。演算子による方法についてさらに詳しくは、GSW\refcite{46}、および双対模型の参考文献Schwarz (1973)\refcite{92}とMandelstam (1974)\refcite{65}を参照されたい。閉弦と開弦の樹木レベルの振幅の関係はKawai, Lewellen, \& Tye (1986)\refcite{54}で与えられている。

\subsection*{第7章}
Mumford (1983)\refcite{73}の2巻は、テータ関数の美しい入門である。1ループの関数は第I巻の前半で論じられている。一般の共形場理論に対するモジュラー不変性からの制約はCardy (1986)\refcite{18}にある。量子場理論の零点エネルギーはColeman \& Weinberg (1973)\refcite{21}で論じられている。宇宙定数問題はWeinberg (1989)\refcite{110}で概説されている。

開超弦のゲージ群$SO(32)$における発散の相殺はGreen \& Schwarz (1985)\refcite{45}で論じられている。境界状態の形式を用いた解析は、Callan, Lovelace, Nappi, \& Yost (1987)\refcite{16}とPolchinski \& Cai (1988)\refcite{84}にある。

\subsection*{第8章}
Kaluza--Klein理論は、Duff, Nilsson, \& Pope (1986)\refcite{30}と、論文再録集Appelquist, Chodos, \& Freund (1987)\refcite{7}で概説されている。Giveon, Porrati, \& Rabinovici (1994)\refcite{42}はトーラス・コンパクト化と$T$双対性を概説している。Polchinski (1997)\refcite{83}は、開弦および向き付けのない弦の$T$双対性と、Dブレーンを概説している。トーラス・コンパクト化についてのさらなる参考文献は、Narain (1986)\refcite{75}とNarain, Sarmadi, \& Witten (1987)\refcite{76}である。

Dixon, Harvey, Vafa, \& Witten (1985)\refcite{29}はオービフォールドの入門である。ねじれた状態の頂点振幅は、Dixon, Friedan, Martinec, \& Shenker (1987)\refcite{28}とHamidi \& Vafa (1987)\refcite{51}で求められている。より広範なオービフォールドの参考文献は第16章のところで挙げる。$c=1$のCFTの空間は、Dijkgraaf, Verlinde, \& Verlinde (1988)\refcite{26}、Ginsparg (1988)\refcite{39}、および総説Ginsparg (1990)\refcite{40}で記述されている。

\subsection*{第9章}
樹木レベルの振幅のユニタリティは、双対模型の総説Schwarz (1973)\refcite{92}とMandelstam (1974)\refcite{65}で論じられている（これらの議論はOCQで行われているのに対し、9.1節ではこの議論のBRST版を示している）。1ループ振幅は、もともとは樹木レベルの振幅にユニタリティを課すことによって構成された。総説としてScherk (1975)\refcite{90}を参照されたい。

高い種数のRiemann面のモジュライ空間は、Martinec (1987)\refcite{68}とD'Hoker \& Phong (1988)\refcite{25}で概説されている。Schottky構成のもう少し詳しい説明は、Mandelstam (1974)\refcite{65}とそこに挙げられた参考文献にある。CFTの縫合せは多くの著者によって論じられてきた。本書の扱いは、Vafa (1987)\refcite{106}、Polchinski (1988)\refcite{81}、Sonoda (1988a,b,1989)\refcite{98}\refcite{99}\refcite{100}、LeClair, Peskin, \& Preitschopf (1989a,b)\refcite{59}\refcite{60}に基づく。開弦への拡張はLewellen (1992)\refcite{61}にある。モジュライ空間のグラフによる分割と最小面積計量の構成はZwiebach (1993)\refcite{113}で概説されている。量子場の摂動論におけるユニタリティの組合せ論についての議論は、't Hooft \& Veltman (1973)\refcite{102}に見いだせる。Fischler--Susskind機構はPolchinski (1988)\refcite{81}で概説されている。光円錐振幅とPolyakov振幅の同値性を示すことに基づく、弦のユニタリティへの別の方法については、D'Hoker \& Giddings (1987)\refcite{23}を参照されたい。光円錐の形式はMandelstam (1973)\refcite{64}とGSW\refcite{46}第11章で論じられている。

本書で記述した形の開弦の場の理論はWitten (1986)\refcite{111}による。弦の場の理論は、Siegel (1988)\refcite{97}、Thorn (1989)\refcite{104}、Zwiebach (1993)\refcite{113}で概説されている。弦の摂動論のBorel総和可能性はGross \& Periwal (1988)\refcite{49}で論じられている。$(2g)!$という振舞いはShenker (1991)\refcite{96}で論じられている。高エネルギーの硬い散乱は、Gross \& Mende (1988)\refcite{48}とGross (1988)\refcite{47}で調べられている。高エネルギー散乱へのもう1つの方法はAmati, Ciafaloni, \& Veneziano (1989)\refcite{6}である。ブーストされた弦の広がりはSusskind (1993)\refcite{101}で記述されている。Hagedorn転移は、Hagedorn (1965)\refcite{50}の最初の研究以来、多くの推測を生んできた。本書の議論は一部Atick \& Witten (1988)\refcite{8}に基づく。

総説Polchinski (1996)\refcite{82}は、本章の主題の多くを扱い、とくに行列模型の入門を与えている。行列模型についてのほかの総説として、Kazakov \& Migdal (1988)\refcite{55}、Klebanov (1992)\refcite{57}、Martinec (1992)\refcite{69}、そしてGinsparg \& Moore (1993)\refcite{41}による広範な総説がある。
